\section{Evaluation}
\subsection{Evaluation Questions}
Our evaluation asks three security-relevant questions. First, does Ghostext recover payloads under fully matched runtime configuration? Second, does Ghostext fail closed under representative mismatch and mutation cases? Third, what practical overhead does the system impose in bits per token and throughput?

\subsection{Setup}
We run Ghostext on a real local backend based on \texttt{llama.cpp} with a Qwen3.5-2B GGUF model. All reported numbers come from rerunnable scripts that emit structured artifacts in \texttt{results/}.

\subsection{Results}
On the real backend, matched-configuration round-trip succeeds on 2/2 cases (100\%). In this run, mean embedding efficiency is 2.199 bits/token, and mean encode throughput is 9.64 tokens/s on the current CPU-only setup. A wrong-prompt sanity check also fails closed with an explicit synchronization error (\texttt{SynchronizationError} in this run).

\begin{table}[t]
\caption{Real backend results (\texttt{llama.cpp} + Qwen3.5-2B GGUF).}
\label{tab:real-backend}
\centering
\begin{tabular}{lrrrr}
\toprule
Case & Success & Packet Bytes & Total Tok & Bits/Tok \\
\midrule
zh\_real\_short & Y & 115 & 441 & 2.086 \\
en\_real\_short & Y & 117 & 405 & 2.311 \\
\bottomrule
\end{tabular}
\end{table}

\textbf{Artifacts.}
Primary outputs are \texttt{results/real\_backend\_eval\_results.json} and \texttt{results/consolidated\_eval\_tables.md}.

\subsection{Interpretation Scope}
These results support recoverability and fail-closed claims under the passive-censor threat model and strict synchronization assumptions. They do not establish robustness under active edits, broad detector resistance, or cross-model/version compatibility.
